% ---
% id: protocol-compiler/thm:choice-encoding
% kind: thm
% title: "Encoding of public choices in the measurement operators"
% summary: >-
%   A public choice of known bias enters as weights inside the measurement
%   operators: the family stays a POVM, fixed weights are reproduced exactly
%   as choice probabilities where the weighted family is emitted, and
%   setting-dependent weights commute with the chain. A fraction outside them
%   is refused.
% prominence: main
% status: implemented
% uses: [def:weighted, lem:weighted, thm:measurements-and-statistics, def:phi, def:compile, def:chaincond, def:spec, def:banda]
% mentions: [rem:relabel, reductions/lem:linear, reductions/rem:intertwining, reductions/ass:tsm, req:no-fixed-fraction, req:setting-only-weighting, req:stochastic-weights, req:weighted-family-emission]
% assumes: [A-COIN, A-LINK]
% refs: []
% provenance: newly-derived
% implements: certiqs_key.protocol_compiler.weights:resolve_povm
% lean:
% notes: "reductions/lem:linear is mentioned, not used: the proof of (vi) writes the linearity argument out and names the lemma only as the general statement (vi) is an instance of. A uses edge would carry the lemma's assumption closure (check_claims --assumptions reductions/lem:linear) into this claim and its dependents."
% ---

The theorem has two halves: (i)--(iv) for a choice of known bias declared on
either party's side, (v)--(viii) for a choice whose weight depends on the
private setting of the apparatus that measures a system, whichever party
draws it, applied after the chain. By (ii) the weight inside
the operators reproduces the declared choice statistics exactly wherever the
weighted family is the one emitted. A measuring party's presentation in front
of a receiver segment that squashes is the other case: there the weighted
family is the physical one handed to the chain, the emitted family is the
chain's, and the first half says what still holds. The
assumption of the output-state test-statistics model that (vii) declares is
\xref{reductions/ass:tsm}; it is named here, and the theorem does not use it.
Remark~\ref{rem:relabel} explains
why the relabelling of (vi) is not a link of the chain; it is the relabelling
of \xref{reductions/rem:intertwining}, which leaves it to the consumer to
apply.

\begin{theorem}
\label{thm:choice-encoding}
Let a per-round public choice of known bias --- a basis choice, a choice
between test and generation rounds, a probe setting --- be declared on either
party's side with alphabet $\Tpub$, weights $(q_t)$ and conditional POVMs
$\{N_{u|t}\}$, under Assumption~\ref{ass:coin}. Then the compiler forms the
choice-weighted family \eqref{eq:weighted}, and (iv) below holds. Where that
family is the emitted one --- the presentation is Alice's, or it is a
measuring party's and the last link of that party's receiver segment
delivers the weighted family itself, on its own alphabet, as its
\tref[target povm]{target POVM}, as on the \tref{Band A} route
(Definition~\ref{def:banda}\,(b)) --- the compiler emits it, and (i)--(iii)
hold as well:
\begin{enumerate}[label=(\roman*),leftmargin=2.2em]
\item the emitted family is a POVM and preserves sector block-diagonality
(Lemma~\ref{lem:weighted}\,(i),(ii)), so the encoding cannot break D(iii) ---
and it is emitted only after the two decisions of
Definition~\ref{def:weighted}, so that what is emitted is binary64 data whose
own predicates D(iii) is then decided on
(Theorem~\ref{thm:measurements-and-statistics}\,(i), and (i$'$) in the
general format);
\item the encoding is \emph{exact}: for every state, the marginal probability
of the choice value $t$ computed from the emitted declaration equals the
declared weight $q_t$ (Lemma~\ref{lem:weighted}\,(iii));
\item the choice value is a component of the joint outcome, hence enters every
downstream quantity only through $P_{\Decl}[\omega]$ and
$\Phimap[\omega]$ of Definition~\ref{def:phi}; no factor $q_t$ is emitted
anywhere outside a POVM element; and
\item the alternative encoding --- declaring the conditional families
$\{N_{u|t}\}$ together with a separate state-independent fraction $q_t$ to be
applied outside the operators --- is never emitted, and a specification
requesting it is refused.
\end{enumerate}
Where the presentation is a measuring party's and the last link of that
party's receiver segment delivers any other family, the weighted family is
the physical receiver family handed to the chain
(Definition~\ref{def:spec}\,(b)) and is not emitted: it is a POVM on the
physical receiver space, and for every unit-trace state there the marginal
probability of $t$ computed from it is $q_t$
(Lemma~\ref{lem:weighted}\,(i), (iii)); no factor $q_t$ is emitted anywhere
outside a POVM element, and (iv) holds. The emitted family is that link's
target POVM, whose own conditions are (C3)'s and check V3's
(Definition~\ref{def:chaincond}), and
nothing here asserts (ii) or (iii) of it; the declared test-statistics model
is the output-state model, as (vii) says of the setting-dependent form.
Let, further, a public choice $f\in\mathcal{F}$ be declared whose weight
depends on the private setting of the apparatus that measures the
\tref{measured system}, as a
stochastic matrix $P(f\,|\,t)$ read through the setting map
$t:\Vset\to\mathcal{T}$ of the family $\{F_v\}_{v\in\Vset}$ the last link of
$\Chrec$ delivered --- the setting-dependent form of
Definition~\ref{def:weighted}; in the general format, at a measured factor,
read through that factor's setting map on the family the last link of that
factor's receiver segment delivered, with $S$ that factor's system --- under
Assumption~\ref{ass:coin}. Then the compiler applies it \emph{after} the
reduction chain, emits the family \eqref{eq:weighted-t},
$N_{v,f}=P(f\,|\,t(v))F_v$ over the support $\Vset_P$, and:
\begin{enumerate}[label=(\roman*),leftmargin=2.2em,resume]
\item the emitted family is a POVM and preserves sector block-diagonality
(Lemma~\ref{lem:weighted}\,(i),(ii)), every element being a scalar multiple
of an element the chain delivered, so the encoding cannot break D(iii); and it
is emitted only after the two decisions of Definition~\ref{def:weighted} on
the products $P(f\,|\,t(v))F_v$, so that what is emitted is binary64 data
whose own predicates D(iii) is then decided on
(Theorem~\ref{thm:measurements-and-statistics}\,(i), and (i$'$) in the
general format);
\item \emph{(soundness: the relabelling commutes with the chain)} for every
channel $\Lambda$ into $S$ and every family $\{G_v\}_{v\in\Vset}$ on its
input space with $\Lambda^{\dagger}(F_v)=G_v$ for every $v$,
\begin{equation}\label{eq:relabel-sound}
  \Lambda^{\dagger}\big(N_{v,f}\big)=P(f\,|\,t(v))\,G_v
  \qquad\text{for every }(v,f)\in\Vset_P :
\end{equation}
the weighted target and the identically weighted input family stand in the
same relation under the \emph{same} channel, with the same target space, the
same sector map, the same constraint entries transported along the chain and
every pull-back unchanged. In this module's setting $G$ is the
$\gamma$-coarse-grained physical family of
Definition~\ref{def:chaincond}\,(C4), $G_v=\sum_{w\in\gamma^{-1}(v)}\Fph_w$,
so the emitted family and the physical family weighted by
$P(f\,|\,t(\gamma(w)))$ satisfy the pullback identity \eqref{eq:pullback}
with the outcome correspondence $(w,f)\mapsto(\gamma(w),f)$ restricted to the
support, wherever (C4) is decidable, and are related as (C4)'s families are
wherever it is not. This is \xref{reductions/lem:linear}\,(i) for the
stochastic matrix from $\Vset$ to $\Vset\times\mathcal{F}$ with entries
$\delta_{vv'}P(f\,|\,t(v'))$;
\item \emph{(what is not asserted)} the exactness clause (ii) is asserted for
the setting-independent form and for it alone. For this form the marginal of
$f$ is the quantity \eqref{eq:fmarginal} of
Lemma~\ref{lem:weighted}\,(iv), a function of the state through the setting
marginal, equal to the declared $q_f$ for every state when $P(f\,|\,t)=q_f$
and state-dependent in general otherwise; nothing here asserts otherwise, no
state-independent fraction is emitted, and the declared test-statistics model
is the output-state model.
What (iv) says of the fixed-fraction encoding holds of this form unchanged;
\item \emph{(refusals)} a weighting presented through anything but the
setting map --- indexed by $\Vset$, by the device outcome, or by any
component of $v$ other than $t(v)$ --- is refused; a $P$ with a negative
entry or a row whose declared-order sum is not exactly $1$ is refused; and
the two emission decisions of Definition~\ref{def:weighted} refuse, for a
product that is not representable or with the predicate that failed. Each
refusal emits no declaration.
\end{enumerate}
\end{theorem}

\begin{proof}
Under the hypothesis of (i)--(iii) the emitted family is the weighted family
itself, on its own alphabet: for Alice's presentation by
Definition~\ref{def:compile}, which emits her measurement slot from the
presentation; for a measuring party's because the emitted family is the
target POVM of the last link (Definition~\ref{def:compile}), which by
hypothesis is the weighted family on $\Vset=\bigsqcup_t\Uset_t\times\{t\}$
--- on the Band A route by Definition~\ref{def:banda}\,(b), with $\gamma$ the
identity on $\Vset=\Vph$ and $F_v=\Fph_v$ entry for entry. So (i) and (ii)
are Lemma~\ref{lem:weighted}, read of the emitted family and against D(iii)
and the emission decisions of Definition~\ref{def:weighted}. For (iii): the
emitted alphabet is $\bigsqcup_t\Uset_t\times\{t\}$ and the emitted matrices
are the $N_{u,t}$; the choice value $t$ is recoverable from the outcome index
alone, and no other emitted slot carries $q_t$. Since $P_{\Decl}[\omega]$ and
$\Phimap[\omega]$ are functions of $\omega$ and the emitted operators only
(Definition~\ref{def:phi}), the claim follows. (iv) is a rule of the
\tref{compilation map}, not a derived statement, and is stated here because it
is the rule the encoding exists to enforce: the emitted declaration carries
the output-state model, and no path of the compiler produces an artefact
carrying a fixed-fraction model. Requirement~\ref{req:no-fixed-fraction}
binds its refusal.

For a measuring party's presentation whose last link delivers any other
family, Definition~\ref{def:spec}\,(b) makes the weighted family the physical
receiver family, the object (C2) and (C4) of Definition~\ref{def:chaincond}
read, and Definition~\ref{def:compile} fills the emitted family from the last
link's target POVM, so the weighted family is not emitted. Lemma~\ref{lem:weighted}\,(i)
and (iii) are stated on any finite-dimensional $S$ and apply with $S$ the
physical receiver space: the weighted family is a POVM there and its choice
marginal is $q_t$ for every unit-trace state there. No slot of the
declaration is filled from the weights: the emitted family is the target
POVM and the remaining slots are the other parts' (Definition~\ref{def:compile}),
so no factor $q_t$ is emitted outside a POVM element, and (iv) is the rule
of the compilation map as before. Nothing is asserted of the emitted family
beyond (C3) and check V3, and nothing more could be in general: (C2) admits
any tabulated total surjection $\gamma$, so $t$ need not be recoverable from
an emitted outcome; and where $\gamma$ keeps the choice value, every emitted
outcome carrying the $t$ of each physical outcome it receives, so that (C4)
summed over the emitted outcomes carrying $t$ reads
$\Lambda^{\dagger}(\Pi_t)=q_t\idon{}$ on the physical receiver space with
$\Pi_t$ the sum of the emitted elements carrying $t$, that fixes $\Pi_t$
only through $\Lambda^{\dagger}$ and does not make it $q_t\idon{\Qsys}$ off
the image of the channel. The
sentence on the test-statistics model is, as in (vii), a statement of what
the emitted declaration carries.

(v) is Lemma~\ref{lem:weighted}\,(i), (ii) for the setting-dependent form,
read against D(iii) and the emission decisions exactly as (i) is.

(vi) By linearity of $\Lambda^{\dagger}$,
$\Lambda^{\dagger}(P(f\,|\,t(v))F_v)=P(f\,|\,t(v))\Lambda^{\dagger}(F_v)
=P(f\,|\,t(v))G_v$, which is \eqref{eq:relabel-sound}; nothing else is used,
and in particular nothing about $\Lambda$ beyond linearity of its adjoint.
That the two weighted families are POVMs is (v) and its mirror on the input
side (the same computation with $G$ for $F$, $\sum_vG_v=\Lambda^{\dagger}
(\sum_vF_v)$ being the identity where $\Lambda^{\dagger}$ is unital, which is
the trace-preservation clause of the obligation of the chain's producer and
not a step of this proof). The identity \eqref{eq:relabel-sound} is the
instance of the post-processing part~(i) of \xref{reductions/lem:linear} for
the relation $G\dsq_{\Lambda}F$ and the stochastic matrix
$\Pi\big((v',f)\,|\,v\big)\coloneqq\delta_{vv'}P(f\,|\,t(v'))$ from $\Vset$ to
$\Vset\times\mathcal{F}$, every column of which sums to
$\sum_fP(f\,|\,t(v))=1$: $(\Pi F)_{(v,f)}=P(f\,|\,t(v))F_v$ and
$(\Pi G)_{(v,f)}=P(f\,|\,t(v))G_v$, and the computation above is that
lemma's argument for $\Pi$, written out. The target space, its
sector map and the transported entries are untouched because a relabelling
changes the index set of the two families and nothing else: every constraint
operator $\hat O_k$ is an operator on $\Qsys$, not a function of $v$, and its
pull-back $\Lambda^{\dagger}(\hat O_k)$ is computed through the same
$\Lambda$. For the compiler's instance of it, substitute (C4)'s
$G_v=\sum_{w\in\gamma^{-1}(v)}\Fph_w$: then
$\Lambda^{\dagger}(N_{v,f})=\sum_{w\in\gamma^{-1}(v)}P(f\,|\,t(\gamma(w)))\Fph_w$,
which is \eqref{eq:pullback} for the family
$\{P(f\,|\,t(\gamma(w)))\Fph_w\}_{(w,f)}$ under the correspondence
$(w,f)\mapsto(\gamma(w),f)$, a total surjection onto $\Vset_P$ because
$\gamma$ is one onto $\Vset$ and the weight depends on $w$ only through
$\gamma(w)$; where (C4) is not decidable from the data,
\eqref{eq:relabel-sound} is still an identity of linear maps, and it says of
the relation the chain's producer asserts and nobody here decides
(Definition~\ref{def:chaincond}\,(C4), second sentence) exactly what it says
of a decided one.

(vii) is Lemma~\ref{lem:weighted}\,(iv): \eqref{eq:fmarginal} is the marginal
of $f$ computed from the emitted declaration, it equals $q_f$ for every state
in the setting-independent case, and it is in general a function of the state
otherwise. The theorem asserts no value for it in that case; the sentence
about the \tref{declared model} is a statement of what the emitted declaration
carries, as (iv) is.

(viii) is the rule of the compilation map that Definition~\ref{def:weighted}
states and Requirements~\ref{req:setting-only-weighting},
\ref{req:stochastic-weights} and~\ref{req:weighted-family-emission} bind,
stated beside (iv) because it is the rule the encoding exists to enforce: a
weighting that reaches an outcome through anything but its setting would make
$f$ a function of the device outcome, which is the inference
Assumption~\ref{ass:coin} excludes; a row of $P$ with non-negative entries
whose sum is not $1$ would make \eqref{eq:forgetf} false at every $v$ of that
setting with $F_v\ne0$, the inner sum over the support \eqref{eq:support}
being that row sum, and a negative entry, which the support never keeps and
which therefore reaches no emitted element, is refused because
Assumption~\ref{ass:coin} reads $P(\,\cdot\,|\,t)$ as the distribution the
choice $f$ is drawn from, and a negative number is no probability; and
the two emission decisions are those of the setting-independent form, taken on
the products of this one.
\end{proof}
